\documentclass[UTF8,a4paper,11pt]{ctexart}
\usepackage[margin=22mm,headheight=14pt,footskip=12mm]{geometry}
\usepackage{amsmath,amssymb,array,longtable,booktabs,tikz,needspace,fancyhdr,hyperref}
\usepackage{microtype}
\IfFontExistsTF{MiSans}{\setCJKmainfont[BoldFont=MiSans Bold]{MiSans}\setCJKsansfont[BoldFont=MiSans Bold]{MiSans}}{}
\hypersetup{colorlinks=false,pdfborder={0 0 0},pdfauthor={Good Learning}}
\newcolumntype{P}[1]{>{\raggedright\arraybackslash}p{#1}}
\setlength{\LTleft}{0pt}\setlength{\LTright}{0pt}
\setlength{\parindent}{0pt}\setlength{\parskip}{.65em}
\linespread{1.18}
\raggedbottom
\setcounter{secnumdepth}{0}
\ctexset{section={format=\Large\bfseries\sffamily,beforeskip=1.35em,afterskip=.65em},
subsection={format=\large\bfseries\sffamily,beforeskip=1em,afterskip=.45em}}
\newcommand{\chaptertitle}[1]{\Needspace{8\baselineskip}\noindent{\LARGE\bfseries\sffamily #1}\par\vspace{4mm}\hrule\vspace{6mm}}
\newcommand{\CompactQuestion}[2]{\par\noindent\textbf{#1.}\ #2\par\vspace{.8em}}
\newcommand{\PracticeQuestion}[3]{\par\Needspace{4\baselineskip}\noindent\begin{minipage}{\linewidth}\textbf{#1.}\ #2\par\noindent\rule{0pt}{#3mm}\par\end{minipage}\par}
\newcommand{\CompactQuestionColon}[2]{\par\noindent\textbf{#1：}\ #2\par\vspace{.8em}}
\newcommand{\PracticeQuestionColon}[3]{\par\Needspace{4\baselineskip}\noindent\begin{minipage}{\linewidth}\textbf{#1：}\ #2\par\noindent\rule{0pt}{#3mm}\par\end{minipage}\par}
\pagestyle{fancy}\fancyhf{}\fancyfoot[C]{\thepage}\renewcommand{\headrulewidth}{0pt}
\begin{document}
\thispagestyle{empty}
\noindent{\sffamily\footnotesize 每日课程}\hfill{\sffamily\footnotesize 学习导图}\par
\vspace{7mm}
\noindent{\sffamily\fontsize{20}{25}\selectfont\bfseries 第 10 天 · 反向传播与 LSTM（宽松排版型）}\par
\vspace{4mm}\noindent\rule{\linewidth}{.85pt}\par\vspace{6mm}
\begin{center}
\begin{tikzpicture}[x=1cm,y=1cm,>=stealth,every node/.style={font=\sffamily}]
\path[use as bounding box] (-8.1,-18.7) rectangle (8.1,.4);
\node[font=\sffamily\Large\bfseries] at (0,-.6) {梯度怎样传回权重和旧状态};
\node[font=\sffamily\small,text=black!55] at (0,-1.25) {先算出每个中间量，再沿路径乘局部导数};
\node[anchor=north west,font=\sffamily\fontsize{22}{24}\selectfont\bfseries,text=black!50] at (-7.35,-2.35) {01};
\node[anchor=north west,text width=5.8cm,align=left,font=\sffamily\large\bfseries] at (-6.199999999999999,-2.3200000000000003) {前向计算};
\node[anchor=north west,text width=5.9cm,align=left,font=\sffamily\small] at (-6.199999999999999,-3.0700000000000003) {逐行乘加，记下激活前后的值\\[.25em]用预测和真实值算损失};
\draw[black!35,line width=.4pt] (-7.35,-5.050000000000001) -- (-0.2999999999999998,-5.050000000000001);
\node[anchor=north west,font=\sffamily\fontsize{22}{24}\selectfont\bfseries,text=black!50] at (0.85,-2.35) {02};
\node[anchor=north west,text width=5.8cm,align=left,font=\sffamily\large\bfseries] at (2.0,-2.3200000000000003) {权重梯度};
\node[anchor=north west,text width=5.9cm,align=left,font=\sffamily\small] at (2.0,-3.0700000000000003) {从损失沿计算路径往回算\\[.25em]一路乘导数，多路加贡献};
\draw[black!35,line width=.4pt] (0.85,-5.050000000000001) -- (7.8999999999999995,-5.050000000000001);
\node[anchor=north west,font=\sffamily\fontsize{22}{24}\selectfont\bfseries,text=black!50] at (0.85,-7.7) {03};
\node[anchor=north west,text width=5.8cm,align=left,font=\sffamily\large\bfseries] at (2.0,-7.67) {跨时间步的梯度};
\node[anchor=north west,text width=5.9cm,align=left,font=\sffamily\small] at (2.0,-8.42) {普通 RNN 反传也会连乘\\[.25em]反复缩小容易导致梯度消失};
\draw[black!35,line width=.4pt] (0.85,-10.4) -- (7.8999999999999995,-10.4);
\node[anchor=north west,font=\sffamily\fontsize{22}{24}\selectfont\bfseries,text=black!50] at (-7.35,-7.7) {04};
\node[anchor=north west,text width=5.8cm,align=left,font=\sffamily\large\bfseries] at (-6.199999999999999,-7.67) {遗忘门与输入门};
\node[anchor=north west,text width=5.9cm,align=left,font=\sffamily\small] at (-6.199999999999999,-8.42) {旧状态乘遗忘门，候选乘输入门\\[.25em]两项相加得到新的细胞状态};
\draw[black!35,line width=.4pt] (-7.35,-10.4) -- (-0.2999999999999998,-10.4);
\node[anchor=north west,font=\sffamily\fontsize{22}{24}\selectfont\bfseries,text=black!50] at (-7.35,-13.05) {05};
\node[anchor=north west,text width=5.8cm,align=left,font=\sffamily\large\bfseries] at (-6.199999999999999,-13.020000000000001) {输出门};
\node[anchor=north west,text width=5.9cm,align=left,font=\sffamily\small] at (-6.199999999999999,-13.770000000000001) {新细胞状态先经 tanh\\[.25em]再乘输出门得到隐藏状态};
\draw[black!35,line width=.4pt] (-7.35,-15.75) -- (-0.2999999999999998,-15.75);
\node[anchor=north west,font=\sffamily\fontsize{22}{24}\selectfont\bfseries,text=black!50] at (0.85,-13.05) {06};
\node[anchor=north west,text width=5.8cm,align=left,font=\sffamily\large\bfseries] at (2.0,-13.020000000000001) {细胞状态直接路径};
\node[anchor=north west,text width=5.9cm,align=left,font=\sffamily\small] at (2.0,-13.770000000000001) {先看单步保留多少旧状态\\[.25em]跨多步时连乘各步遗忘门};
\draw[black!35,line width=.4pt] (0.85,-15.75) -- (7.8999999999999995,-15.75);
\draw[->,black!55,line width=.8pt] (-.12,-3.7) -- (.55,-3.7);
\draw[->,black!55,line width=.8pt] (7.85,-5.45) -- (7.85,-7.2);
\draw[->,black!55,line width=.8pt] (.52,-9.05) -- (-.15,-9.05);
\draw[->,black!55,line width=.8pt] (-7.85,-10.8) -- (-7.85,-12.55);
\draw[->,black!55,line width=.8pt] (-.12,-14.4) -- (.55,-14.4);
\node[font=\sffamily\small,text=black!55] at (0,-18.1) {10 分钟热身 · 25 分钟讲解 · 40 分钟练习 · 15 分钟订正};
\end{tikzpicture}
\end{center}
\clearpage

\chaptertitle{第 10 天 · 反向传播与 LSTM}
\textbf{2026 年 11 月 24 日｜K11、K12｜90 分钟}

今天要交两样结果。一份权重梯度的计算过程，一张标出遗忘门、输入门和输出门的 LSTM 图。先复习上一课的两层前向。算反向传播需要会求平方函数的导数，并记得把 ReLU 前的 $z$ 和处理后的 $a$ 分开写。LSTM 练习会给定门值和所需的 $\tanh$ 近似值。

\section{90 分钟安排}
\begin{itemize}
\item 10 分钟 · 闭卷做 H1，按权重矩阵逐行写乘加，标出各量形状。
\item 25 分钟 · 看两个单元的例题，再做各自的填步与变式。一个单元算权重梯度，一个单元画 LSTM 的信息流。
\item 40 分钟 · 做 F1、M1、M2。建议填空用 5 分钟，反传和 LSTM 分别用 20 分钟和 15 分钟。
\item 15 分钟 · 看后面的解析。合上解析后再算一个权重的梯度，并重新画出三个门。
\end{itemize}
H1 是错题本第 21 条的复做题。其余题目均为新编，F1 对应原计划的填空考点，M1 和 M2 用于原计划的计算练习，分别练反向传播与序列状态。工作区可核对原计划的题号，讲义包含做题所需的全部信息。

\subsection{举一反三 · 先复做前向传播}
\PracticeQuestion{H1}{输入 $x=(2,-1)^\top$。第一层权重的两行是 $(1,-1)$ 和 $(-2,1)$，偏置 $b^{(1)}=(0,1)^\top$，使用 ReLU。第二层权重为 $(3,2)$，偏置为 $-1$，使用恒等函数。逐行写乘加，求第一层的 $z^{(1)}$、$a^{(1)}$ 和第二层的输出。写出两层各自的输入、权重、偏置和输出形状。注明被置零的神经元在哪一层、排第几个。}{55}
\clearpage
\section{单元 1 · 从损失往回算权重梯度}
前向传播用权重算预测，再用预测和真实值算损失。反向传播沿同一条计算路径往回走，求每个权重微小变化时，损失会往哪个方向变。

以 $z=wx+b$、$a=\operatorname{ReLU}(z)$、$\hat y=va+c$、$L=\tfrac12(\hat y-y)^2$ 为例。链式法则沿路径把相邻量之间的导数相乘。如果变化沿多条路径抵达同一个量，就把各路径的贡献相加。

\[\frac{\partial L}{\partial w}=\frac{\partial L}{\partial\hat y}\frac{\partial\hat y}{\partial a}\frac{\partial a}{\partial z}\frac{\partial z}{\partial w}=(\hat y-y)\,v\,\operatorname{ReLU}'(z)\,x.\]
ReLU 在 $z>0$ 时斜率为 1，在 $z<0$ 时斜率为 0。在 $z=0$ 处不可导，做题按题设约定，本课计算避开该点。算出的梯度可用于下一步更新参数，例如 $w_{\rm new}=w-\eta\,\partial L/\partial w$。

\subsection{完整例题 · 哪个权重使损失变大}
令 $x=2$，隐藏层 $w=1,b=-1$，输出层 $v=3,c=0$，真实值 $y=1$。

\begin{enumerate}
\item 前向先算 $z=1(2)-1=1$、$a=1$、$\hat y=3(1)=3$。记下 $z>0$，反传会用到它。
\item 从损失开始。因 $L=\tfrac12(\hat y-y)^2$，对预测求导得 $\partial L/\partial\hat y=3-1=2$。
\item 求 $v$ 的梯度。输出 $\hat y=va+c$ 对 $v$ 的导数是 $a=1$，所以 $\partial L/\partial v=2(1)=2$。
\item 求 $w$ 的梯度。输出对 $a$ 的导数是 3，ReLU 的斜率是 1，$z$ 对 $w$ 的导数是输入 2。四项相乘得 $2(3)(1)(2)=12$。验算方向时，想象稍微增大 $w$：预测 3 会继续增大，离真实值 1 更远，梯度为正合理。
\end{enumerate}
\subsection{举一反三}
\PracticeQuestion{A1}{沿用例题参数，把真实值改成 $y=2$。已知前向结果未变，$\partial L/\partial\hat y=1$。补出 $\partial L/\partial w$ 并标出 ReLU 的斜率。}{23}
\PracticeQuestion{A2}{恢复 $y=1$，把输入换成 $x=-2$。重新算 $z,a,\hat y$，再求 $\partial L/\partial w$、$\partial L/\partial v$ 和输出偏置梯度 $\partial L/\partial c$。指出相较例题哪一条局部导数变为 0。}{35}
\clearpage
\section{单元 2 · LSTM 怎样保留、写入和输出}
普通 RNN 反复用上一步的隐藏状态 $h_{t-1}$ 和当前输入 $x_t$ 计算新的 $h_t$。跨很多步反传时会连乘许多局部导数。如果梯度反复被缩小，较早的时间步就很难收到训练信号，这叫梯度消失。例如，连续乘 10 次 $0.5$ 只剩约 $0.001$。Sigmoid 在两端趋于平坦，导数接近 0，也会缩小传回去的梯度。

LSTM 额外维护细胞状态 $c_t$。给定 $x_t$ 和 $h_{t-1}$，三个 Sigmoid 门给出介于 0 和 1 之间的系数。候选内容 $\tilde c_t$ 通常经过 $\tanh$，可以是负数。以下用标量演示。向量版的门值与状态对应分量逐项相乘。

\begin{itemize}
\item 遗忘门 $f_t$ 乘旧状态 $c_{t-1}$，决定保留多少旧信息。
\item 输入门 $i_t$ 乘候选内容 $\tilde c_t$，决定写入多少新信息。
\item 输出门 $o_t$ 乘 $\tanh(c_t)$，决定当前隐藏状态 $h_t$ 输出多少。
\end{itemize}
\[c_t=f_t c_{t-1}+i_t\tilde c_t,\qquad h_t=o_t\tanh(c_t).\]
画图时，把「旧状态乘遗忘门」与「候选乘输入门」的两条支路接到一个加号，得到 $c_t$。从 $c_t$ 另接一条线，经 $\tanh$ 后乘输出门，再得到 $h_t$。

\subsection{完整例题 · 手算一个时间步}
旧状态 $c_{t-1}=2$，门值 $f_t=0.75$、$i_t=0.5$、$o_t=0.4$，候选 $\tilde c_t=-1/2$。已知 $\tanh(1.25)\approx0.848$。

\begin{enumerate}
\item 旧状态通过遗忘门，保留 $0.75(2)=1.5$。
\item 候选通过输入门，写入 $0.5(-1/2)=-0.25$。候选可以为负，所以写入后状态可能变小。
\item 加起来得 $c_t=1.5-0.25=1.25$。下一时间步的细胞状态来自这里。
\item 对新状态取 $\tanh$ 后乘输出门，得到 $h_t\approx0.4(0.848)=0.3392$。用同一个 $c_t$ 检查，减小输出门会改变 $h_t$，当前的 $c_t$ 仍为 1.25。
\end{enumerate}
只看 $c_{t-1}$ 到 $c_t$ 的直接路径，并把门值和候选视为给定量，局部梯度等于 $f_t$。遗忘门接近 1 时，梯度沿这条路径衰减较慢。完整网络还包含经过隐藏状态和门值的其他路径。

\subsection{举一反三}
\PracticeQuestion{B1}{沿用例题，已知保留量为 1.5，写入量为 $-0.25$。把输出门改为 $0.2$，补出 $c_t$ 和 $h_t$。}{13}
\PracticeQuestion{B2}{沿用例题，把遗忘门改为 $0.25$，候选改为 $+1/2$。给定 $\tanh(0.75)\approx0.635$，重新计算保留量、写入量、$c_t$ 和 $h_t$，指出变化发生在哪两条支路。}{23}
\clearpage
\section{独立练习}
先写计算过程，再翻答案。近似值保留三位小数。梯度的形状要和对应权重的形状一致。

\PracticeQuestion{F1}{LSTM 用 $\underline{\qquad}$ 门控制旧细胞状态保留多少，用 $\underline{\qquad}$ 门控制候选内容写入多少，用 $\underline{\qquad}$ 门控制 $\tanh(c_t)$ 输出多少。常用 $\underline{\qquad}$ 函数把这些门的值约束在 0 到 1 之间。}{12}
\PracticeQuestion{M1}{输入为 $(1,2)^\top$。隐藏层权重两行分别是 $(1,1)$、$(2,-2)$，偏置 $(0,1)^\top$，用 ReLU。输出层权重 $v=(2,-1)$，偏置 $c=1$，用恒等函数。真值 $y=5$，损失 $L=\tfrac12(\hat y-y)^2$。先写 $z^{(1)}$、$a^{(1)}$、$\hat y$ 和 $L$。再分别求输出层权重的两个梯度、隐藏层权重矩阵的梯度和两层的偏置梯度，标出权重梯度的形状。解释哪一行隐藏层梯度为零。}{63}
\PracticeQuestion{M2}{标量 LSTM 初始 $c_0=2$。第 1 步的 $(f_1,i_1,\tilde c_1,o_1)=(0.5,0.5,0.5,0.5)$，第 2 步的 $(f_2,i_2,\tilde c_2,o_2)=(0.8,0.5,-0.4,0.25)$。给定 $\tanh(1.25)\approx0.848$、$\tanh(0.8)\approx0.664$。分别计算两个时间步的 $c_t$ 和 $h_t$。画出第 2 步的三个门并标注相乘对象。若只沿给定门值下的细胞状态直接路径反传，求 $\partial c_2/\partial c_0$。}{63}
\clearpage
\section{核对答案与解析}
先自己算完。如果一个梯度与例题不符，从前向时记录的正负号开始查。

\subsection{热身、填步与变式}
\textbf{H1.} 第 1 层两行是 $1(2)+(-1)(-1)+0=3$、$(-2)(2)+1(-1)+1=-4$，所以 $z^{(1)}=(3,-4)^\top$、$a^{(1)}=(3,0)^\top$。隐藏层第 2 个神经元被置零。第 2 层输出为 $3(3)+2(0)-1=8$。两层的输入都是 $2\times1$。第 1 层权重为 $2\times2$，偏置和输出为 $2\times1$。第 2 层权重为 $1\times2$，偏置和输出为 $1\times1$。

\textbf{A1.} $z=1>0$，ReLU 斜率为 1。$\partial L/\partial w=1(3)(1)(2)=6$。

\textbf{A2.} $z=1(-2)-1=-3$、$a=0$、$\hat y=0$。损失对预测的导数为 $-1$。ReLU 的斜率为 0，所以 $\partial L/\partial w=(-1)(3)(0)(-2)=0$。输出权重梯度为 $(-1)a=0$。输出偏置梯度为 $(-1)(1)=-1$，偏置仍直接影响预测。

\textbf{B1.} $c_t=1.5-0.25=1.25$，$h_t\approx0.2(0.848)=0.170$。当前步的输出门只在计算 $h_t$ 时用到。

\textbf{B2.} 遗忘门支路是 $0.25(2)=0.5$，输入门支路是 $0.5(1/2)=0.25$。相加得 $c_t=0.75$。输出门仍是 $0.4$，所以 $h_t\approx0.4(0.635)=0.254$。

\clearpage
\subsection{混合题答案}
\textbf{F1.} 遗忘门、输入门、输出门、Sigmoid。三个门分别乘旧状态、候选内容和经过 $\tanh$ 的新状态。

\textbf{M1.} 先算前向。第一行是 $1+2+0=3$，第二行是 $2-4+1=-1$。因此 $z^{(1)}=(3,-1)^\top$、$a^{(1)}=(3,0)^\top$。$\hat y=2(3)+(-1)(0)+1=7$，$L=\tfrac12(7-5)^2=2$。

输出层收到梯度 $\hat y-y=2$。每个输出权重乘的是对应的隐藏激活，权重梯度为 $(2\cdot3,2\cdot0)=(6,0)$，形状 $1\times2$。输出偏置梯度为 $2$。

反传到两个隐藏激活的梯度分别为 $2(2)=4$、$2(-1)=-2$。逐项乘 ReLU 的斜率 $(1,0)^\top$ 后得到 $(4,0)^\top$，也就是隐藏偏置梯度。第 1 行的两个权重分别乘输入 1、2，得 $(4,8)$。第 2 行得 $(0,0)$，因为对应的 $z=-1$，ReLU 斜率为 0。隐藏权重的梯度为

\[\frac{\partial L}{\partial W^{(1)}}=\begin{pmatrix}4&8\\0&0\end{pmatrix},\qquad \text{形状为 }2\times2.\]
\textbf{M2.} 第 1 步保留 $0.5(2)=1$，写入 $0.5(0.5)=0.25$，所以 $c_1=1.25$、$h_1\approx0.5(0.848)=0.424$。第 2 步保留 $0.8(1.25)=1$，写入 $0.5(-0.4)=-0.2$，所以 $c_2=0.8$、$h_2\approx0.25(0.664)=0.166$。

第 2 步图的三条线可写为

\[\begin{aligned}
c_1&\xrightarrow{\ \times f_2\ }\ f_2 c_1\\
\tilde c_2&\xrightarrow{\ \times i_2\ }\ i_2\tilde c_2\\
f_2 c_1+i_2\tilde c_2&=c_2\xrightarrow{\ \tanh\ }\tanh(c_2)\xrightarrow{\ \times o_2\ } h_2.
\end{aligned}\]
按题意只看直接路径，$c_0$ 到 $c_1$ 乘 $f_1$，$c_1$ 到 $c_2$ 乘 $f_2$。两步连乘得 $0.5(0.8)=0.4$。

\section{学完后自检}
从 M1 任选一个隐藏权重，闭卷写出从损失到它的四个导数。再从 M2 画出旧状态、候选内容和新状态各接哪一个门。发 H1 的逐行算式、M1 的梯度矩阵和 M2 的门控图或作答照片，并注明用过提示的位置。

\end{document}
